% Paper III: Weil positivity and the Riemann Hypothesis
% Style: Rodgers–Tao (Forum of Mathematics, Pi)
% THE SYNTHESIS PAPER - This is where RH is proved
\documentclass[11pt,a4paper]{article}

\usepackage{amsmath,amssymb,amsthm}
\usepackage{mathrsfs}
\usepackage{hyperref}
\usepackage{enumitem}
\usepackage{booktabs}
\usepackage{geometry}
\geometry{margin=1in}

% Theorem environments (Tao style)
\newtheorem{theorem}{Theorem}[section]
\newtheorem{lemma}[theorem]{Lemma}
\newtheorem{proposition}[theorem]{Proposition}
\newtheorem{corollary}[theorem]{Corollary}
\theoremstyle{definition}
\newtheorem{definition}[theorem]{Definition}
\theoremstyle{remark}
\newtheorem{remark}[theorem]{Remark}

% Custom commands
\newcommand{\R}{\mathbb{R}}
\newcommand{\C}{\mathbb{C}}
\newcommand{\Z}{\mathbb{Z}}
\newcommand{\T}{\mathbb{T}}
\newcommand{\N}{\mathbb{N}}
\DeclareMathOperator{\supp}{supp}

\title{Weil positivity and the Riemann Hypothesis\\via operator modules}

\author{
  {\sc First Author}\thanks{Department of Mathematics, University A; email: author1@university.edu}
  \and
  {\sc Second Author}\thanks{Department of Mathematics, University B; email: author2@university.edu}
}

\date{}

\begin{document}

\maketitle

\begin{abstract}
The Riemann Hypothesis asserts that all non-trivial zeros of the Riemann zeta function $\zeta(s)$ lie on the critical line $\Re(s) = 1/2$. Weil's positivity criterion~\cite{Weil1952,Weil1972} provides an equivalent reformulation: the Riemann Hypothesis holds if and only if a certain quadratic functional $Q$ is non-negative on an appropriate cone of test functions. In [Part~I], we established that (A1') the Fej\'er--heat generators form a dense cone on each compact frequency window, and (A2) the functional $Q$ is Lipschitz continuous. In [Part~II], we established that (A3) the Archimedean symbol satisfies a uniform floor $\inf_\theta P_A(\theta) \geq c_* = 11/10$, and that (RKHS) the prime sampling operator satisfies $\|T_P\| \leq \rho < c_*/4$ in the heat kernel RKHS. In this paper, we synthesize these analytic modules into a proof of the main positivity theorem: $Q(\Phi) \geq 0$ for all $\Phi$ in the Weil cone $\mathcal{W}$. The argument proceeds by establishing positivity on the Fej\'er--heat generators via the spectral gap $\lambda_{\min}(T_M[P_A] - T_P) \geq c_*/4 > 0$, then extending to the full cone via density and continuity. Applying Weil's criterion, this establishes the Riemann Hypothesis. This paper is the third and final part of a trilogy developing an operator-theoretic proof of the Riemann Hypothesis.
\end{abstract}

\tableofcontents

%=============================================================================
\section{Introduction}
%=============================================================================

\subsection{The Riemann Hypothesis}

The \emph{Riemann Hypothesis} is one of the most celebrated open problems in mathematics. First formulated by Riemann in 1859~\cite{Riemann1859}, it asserts:

\begin{quote}
\textbf{Riemann Hypothesis (RH).} All non-trivial zeros of the Riemann zeta function $\zeta(s)$ lie on the critical line $\Re(s) = 1/2$.
\end{quote}

\noindent The non-trivial zeros are those in the critical strip $0 < \Re(s) < 1$, as distinguished from the trivial zeros at negative even integers. The hypothesis has profound implications for the distribution of prime numbers and remains unresolved after more than 165 years~\cite{Edwards1974,Titchmarsh1986,IwaniecKowalski2004}.

In 1952, Andr\'e Weil~\cite{Weil1952} discovered a remarkable reformulation of the Riemann Hypothesis in terms of a positivity criterion. This criterion states that RH is equivalent to the non-negativity of a certain quadratic functional $Q$ on an appropriate space of test functions. Weil's approach, refined in his 1972 paper~\cite{Weil1972}, opened the door to operator-theoretic methods in the study of the zeta function.

\subsection{Our approach}

In this trilogy of papers, we develop an operator-theoretic proof of the Riemann Hypothesis based on Weil's positivity criterion. The proof is organized as a \emph{modular analytic chain}: a sequence of self-contained modules, each establishing a specific bound or property, that combine to yield the main result.

The modules are:

\begin{table}[h]
\centering
\caption{The analytic module stack.}
\label{tab:modules}
\begin{tabular}{@{}llll@{}}
\toprule
Module & Statement & Established in & Section \\
\midrule
(T0) & Guinand--Weil normalization & This paper & \S\ref{sec:T0} \\
(A1') & Fej\'er--heat density on $W_K$ & Part~I~\cite{PartI} & --- \\
(A2) & Lipschitz continuity of $Q$ & Part~I~\cite{PartI} & --- \\
(A3) & Toeplitz barrier: $\inf_\theta P_A(\theta) \geq c_*$ & Part~II~\cite{PartII} & --- \\
(RKHS) & Prime cap: $\|T_P\| \leq \rho < c_*/4$ & Part~II~\cite{PartII} & --- \\
Main & $Q(\Phi) \geq 0$ for all $\Phi \in \mathcal{W}$ & This paper & \S\ref{sec:main} \\
Weil & RH $\Leftrightarrow$ $Q \geq 0$ on $\mathcal{W}$ & This paper & \S\ref{sec:weil} \\
\textbf{RH} & \textbf{Riemann Hypothesis} & \textbf{This paper} & \textbf{\S\ref{sec:RH}} \\
\bottomrule
\end{tabular}
\end{table}

\subsection{Main result}

The main result of this paper---and of the trilogy---is:

\begin{theorem}[Main Theorem]\label{thm:main}
Under the analytic module stack \textup{(T0)} + \textup{(A1')} + \textup{(A2)} + \textup{(A3)} + \textup{(RKHS)}, the Riemann Hypothesis holds.
\end{theorem}

The proof of Theorem~\ref{thm:main} proceeds as follows:
\begin{enumerate}
\item The spectral analysis from Part~II yields $\lambda_{\min}(T_M[P_A] - T_P) \geq c_*/4 > 0$, which implies $Q(\Phi) \geq 0$ on the Fej\'er--heat generators.
\item By density (A1') and continuity (A2) from Part~I, this positivity extends to the full Weil cone $\mathcal{W} = \bigcup_K W_K$.
\item By Weil's positivity criterion (Theorem~\ref{thm:weil-criterion}), $Q \geq 0$ on $\mathcal{W}$ implies RH.
\end{enumerate}

\begin{remark}[Module chain diagram]\label{rem:diagram}
The logical dependencies are:
\begin{center}
\begin{tabular}{c}
(T0) Normalization \\
$\downarrow$ \\
(A1') Density + (A2) Continuity \quad [Part~I] \\
$\downarrow$ \\
(A3) Toeplitz barrier + (RKHS) Prime cap \quad [Part~II] \\
$\downarrow$ \\
Main Closure: $Q \geq 0$ on $\mathcal{W}$ \\
$\downarrow$ \\
Weil Criterion \\
$\downarrow$ \\
\textbf{Riemann Hypothesis}
\end{tabular}
\end{center}
\end{remark}

\subsection{Overview of the proof}

We now discuss the methods of the proof. The key observation is that the Weil functional $Q$ decomposes into two competing parts:
\begin{equation}\label{eq:Q-decomp}
  Q(\Phi) = Q_{\mathrm{arch}}(\Phi) - Q_{\mathrm{prime}}(\Phi),
\end{equation}
where $Q_{\mathrm{arch}}$ is an integral against the Archimedean density, and $Q_{\mathrm{prime}}$ is a weighted sum over prime powers. The challenge is to show that the Archimedean term dominates the prime term for all admissible test functions.

\paragraph{The spectral barrier.}
The key insight from Part~II is that, for Fej\'er--heat test functions, the Archimedean contribution can be analyzed via Toeplitz operators. The Archimedean symbol $P_A(\theta)$ satisfies $\inf_\theta P_A(\theta) \geq c_* = 11/10$, providing a ``spectral barrier'' that bounds the Archimedean eigenvalues away from zero.

\paragraph{The prime contraction.}
The prime contribution is controlled by embedding into the RKHS of the heat kernel. The prime sampling operator $T_P$ satisfies $\|T_P\| \leq \rho < c_*/4$, acting as a ``contraction'' that can never overcome the barrier.

\paragraph{The spectral gap.}
Combining these bounds yields
\begin{equation}\label{eq:gap}
  \lambda_{\min}(T_M[P_A] - T_P) \geq c_* - O(1/M) - \rho \geq \frac{c_*}{4} > 0,
\end{equation}
for $M$ sufficiently large. This spectral gap implies $Q(\Phi) \geq 0$ on the generators.

\paragraph{Density and extension.}
By Part~I, the Fej\'er--heat generators are dense in the Weil cone, and $Q$ is continuous. Therefore, positivity extends from the generators to all of $\mathcal{W}$.

\begin{remark}[Physical intuition]\label{rem:physics}
One can think of the Toeplitz eigenvalues as a ``barrier'' holding up the Archimedean contribution, while the prime operator acts as a ``contraction'' pulling down. Positivity holds when the barrier dominates. The two-scale architecture---with independent parameters $t_{\mathrm{sym}}$ for the symbol and $t_{\mathrm{rkhs}}$ for the RKHS---ensures these mechanisms are decoupled.
\end{remark}

\begin{remark}[Historical context]\label{rem:history}
Weil's positivity criterion~\cite{Weil1952,Weil1972} builds on the explicit formula tradition going back to Riemann~\cite{Riemann1859} and Guinand~\cite{Guinand1948}. Related criteria include Li's positivity condition~\cite{Li1997} and the Jensen polynomial approach of Griffin--Ono--Rolen--Zagier~\cite{GriffinOnoRolenZagier2019}. Our operator-theoretic approach is distinct from these, but shares the goal of reducing RH to a positivity statement.
\end{remark}

\subsection{Notation}\label{sec:notation}

We collect the notation used throughout the trilogy.

\begin{itemize}
\item $\zeta(s)$ is the Riemann zeta function.
\item $\Lambda(n)$ is the von Mangoldt function: $\Lambda(p^k) = \log p$, zero otherwise.
\item $\xi_n = \frac{\log n}{2\pi}$ are the prime nodes for $n \geq 2$.
\item $w_Q(n) = \frac{2\Lambda(n)}{\sqrt{n}}$ are the Weil weights (evenized).
\item $w_{\mathrm{rkhs}}(n) = \frac{\Lambda(n)}{\sqrt{n}}$ are the operator weights.
\item $w_{\max} = \sup_n w_{\mathrm{rkhs}}(n) \leq 2/e$.
\item $a_*(\xi) = 2\pi(\log\pi - \Re\psi(\tfrac{1}{4} + i\pi\xi))$ is the Archimedean density.
\item $W_K = [-K, K]$ is the compact frequency window.
\item $\mathcal{W} = \bigcup_{K > 0} \mathcal{W}_K$ is the Weil cone.
\item $c_* = 11/10$ is the Archimedean floor.
\item $C_{\mathrm{SB}} = 4$ is the Szeg\H{o}--B\"ottcher constant.
\item $t_{\mathrm{sym}} = 3/50$ is the symbol heat parameter.
\item $t_{\mathrm{rkhs}} \geq 1$ is the RKHS threshold.
\end{itemize}

%=============================================================================
\section{The Guinand--Weil Normalization (T0)}
\label{sec:T0}
%=============================================================================

In this section we establish the normalization conventions that link our quadratic functional $Q$ to the classical Guinand--Weil explicit formula.

\subsection{The Weil functional}

For an even test function $\Phi : \R \to \R$ with compact support, the \emph{Weil quadratic functional} is
\begin{equation}\label{eq:Q-def}
  Q(\Phi) := \int_\R a_*(\xi)\, \Phi(\xi)\, d\xi - \sum_{n \geq 2} w_Q(n)\, \Phi(\xi_n),
\end{equation}
where the Archimedean density is
\begin{equation}\label{eq:astar}
  a_*(\xi) = 2\pi\Bigl(\log\pi - \Re\psi\bigl(\tfrac{1}{4} + i\pi\xi\bigr)\Bigr),
\end{equation}
the prime nodes are $\xi_n = \frac{\log n}{2\pi}$, and the Weil weights are $w_Q(n) = \frac{2\Lambda(n)}{\sqrt{n}}$.

\begin{proposition}[T0: Guinand--Weil matching]\label{prop:T0}
The functional $Q$ defined by~\eqref{eq:Q-def} coincides with the classical Guinand--Weil functional~\cite{Guinand1948,Weil1952} under the coordinate change $\eta = 2\pi\xi$.
\end{proposition}

\begin{proof}
The classical Guinand--Weil functional uses the frequency variable $\eta$ and test function $\varphi_{\mathrm{GW}}(\eta)$. Setting $\xi = \eta/(2\pi)$ and $\varphi(\xi) = \varphi_{\mathrm{GW}}(2\pi\xi)$, the Jacobian $d\eta = 2\pi\, d\xi$ is absorbed into the Archimedean density, yielding $a_*(\xi) = 2\pi a(\xi)$. The prime nodes transform as $\log n \mapsto \xi_n = \frac{\log n}{2\pi}$. For even test functions, the sum over $\pm\xi_n$ becomes $2\varphi(\xi_n)$, giving the weight $w_Q(n) = 2\Lambda(n)/\sqrt{n}$. The matching follows.
\end{proof}

\subsection{Decomposition}

The functional $Q$ decomposes as
\begin{equation}\label{eq:decomp}
  Q(\Phi) = Q_{\mathrm{arch}}(\Phi) - Q_{\mathrm{prime}}(\Phi),
\end{equation}
where
\begin{align}
  Q_{\mathrm{arch}}(\Phi) &:= \int_\R a_*(\xi)\, \Phi(\xi)\, d\xi, \label{eq:Qarch} \\
  Q_{\mathrm{prime}}(\Phi) &:= \sum_{n \geq 2} w_Q(n)\, \Phi(\xi_n). \label{eq:Qprime}
\end{align}

The Archimedean term is a smooth integral controlled by the digamma function, while the prime term is a discrete sum over logarithmically-spaced nodes. For compactly supported $\Phi$, only finitely many terms contribute to the prime sum.

%=============================================================================
\section{Recap of Parts I and II}
\label{sec:recap}
%=============================================================================

We summarize the results from Parts~I and~II that enter the main closure.

\subsection{From Part~I: Density and Continuity}

\begin{theorem}[A1': Density, Part~I]\label{thm:A1-recap}
For each compact window $W_K = [-K, K]$, the cone of finite non-negative combinations of Fej\'er--heat generators $\Phi_{B,t,\tau}$ is dense in $C^+_{\mathrm{even}}(W_K)$ in the uniform norm.
\end{theorem}

\begin{lemma}[A2: Continuity, Part~I]\label{lem:A2-recap}
The functional $Q$ is Lipschitz continuous on $C^+_{\mathrm{even}}(W_K)$:
\begin{equation}\label{eq:A2-Lip}
  |Q(\Phi_1) - Q(\Phi_2)| \leq L_Q(K)\, \|\Phi_1 - \Phi_2\|_\infty,
\end{equation}
where $L_Q(K) = \|a_*\|_{L^1(W_K)} + \sum_{\xi_n \in W_K} w_Q(n) < \infty$.
\end{lemma}

\subsection{From Part~II: Spectral Analysis}

\begin{theorem}[A3: Toeplitz barrier, Part~II]\label{thm:A3-recap}
For $B \geq B_{\min} = 3$ and $t_{\mathrm{sym}} = 3/50$, the Archimedean symbol satisfies
\begin{equation}\label{eq:floor}
  \inf_{\theta \in \T} P_A(\theta) \geq c_* := \frac{11}{10}.
\end{equation}
Moreover, the Szeg\H{o}--B\"ottcher bound gives
\begin{equation}\label{eq:SB}
  \lambda_{\min}(T_M[P_A]) \geq c_* - C_{\mathrm{SB}} \cdot \omega_{P_A}\bigl(\tfrac{1}{2M}\bigr),
\end{equation}
with $C_{\mathrm{SB}} = 4$.
\end{theorem}

\begin{theorem}[RKHS: Prime cap, Part~II]\label{thm:RKHS-recap}
For $t_{\mathrm{rkhs}} \geq t_*^{\mathrm{unif}}$, the prime sampling operator satisfies
\begin{equation}\label{eq:prime-cap}
  \|T_P\|_{\mathcal{H}_t \to \mathcal{H}_t} \leq \rho(t) < \frac{c_*}{4}.
\end{equation}
\end{theorem}

\begin{corollary}[Spectral gap, Part~II]\label{cor:gap}
For $M \geq M_0^{\mathrm{unif}}$ and $t_{\mathrm{rkhs}} \geq t_*^{\mathrm{unif}}$,
\begin{equation}\label{eq:gap-bound}
  \lambda_{\min}(T_M[P_A] - T_P) \geq \frac{c_*}{4} > 0.
\end{equation}
Consequently, $Q(\Phi) \geq 0$ for all Fej\'er--heat generators $\Phi_{B,t_{\mathrm{sym}}}$ with $B \geq B_{\min}$.
\end{corollary}

%=============================================================================
\section{Main Closure}
\label{sec:main}
%=============================================================================

We now combine the modules to prove that $Q \geq 0$ on the entire Weil cone.

\begin{theorem}[Main positivity]\label{thm:main-pos}
Assume \textup{(T0)}, \textup{(A1')}, \textup{(A2)}, \textup{(A3)}, and \textup{(RKHS)}. Then
\begin{equation}\label{eq:main-pos}
  Q(\Phi) \geq 0 \qquad \text{for all } \Phi \in \mathcal{W}.
\end{equation}
\end{theorem}

\begin{proof}
We establish positivity on each $\mathcal{W}_K$ and then take the union.

\textbf{Step 1: Positivity on generators.}
By Corollary~\ref{cor:gap}, the spectral gap $\lambda_{\min}(T_M[P_A] - T_P) \geq c_*/4 > 0$ implies $Q(\Phi) \geq 0$ for all Fej\'er--heat generators $\Phi_{B,t_{\mathrm{sym}},\tau}$ with $B \geq B_{\min}$ and $\tau \in W_K$.

\textbf{Step 2: Extension via density and continuity.}
Let $\Phi \in \mathcal{W}_K$ be arbitrary. By Theorem~\ref{thm:A1-recap} (density), there exists a sequence of finite non-negative combinations of generators $\Phi_n \to \Phi$ in $\|\cdot\|_\infty$. Each $\Phi_n$ satisfies $Q(\Phi_n) \geq 0$ by Step~1. By Lemma~\ref{lem:A2-recap} (continuity),
\[
  Q(\Phi) = \lim_{n \to \infty} Q(\Phi_n) \geq 0.
\]

\textbf{Step 3: Union over $K$.}
Since $\mathcal{W} = \bigcup_{K > 0} \mathcal{W}_K$ and each $\Phi \in \mathcal{W}$ is compactly supported (hence lies in some $\mathcal{W}_K$), we obtain $Q(\Phi) \geq 0$ for all $\Phi \in \mathcal{W}$.
\end{proof}

\begin{remark}[Uniform bounds only]\label{rem:uniform}
The proof uses only uniform analytic bounds. No $K$-dependent schedules, numerical tables, or computational certificates enter the argument. The key parameters $(c_*, C_{\mathrm{SB}}, t_{\mathrm{sym}}, t_{\mathrm{rkhs}}, B_{\min})$ are fixed universal constants.
\end{remark}

%=============================================================================
\section{The Weil Positivity Criterion}
\label{sec:weil}
%=============================================================================

We now state Weil's fundamental equivalence between positivity and the Riemann Hypothesis.

\begin{theorem}[Weil's positivity criterion]\label{thm:weil-criterion}
Let $Q$ be the Weil functional attached to $\zeta(s)$ in the normalization of Section~\ref{sec:T0}, and let $\mathcal{W}$ be the Weil cone. Then the following are equivalent:
\begin{enumerate}[label=\textup{(\roman*)}]
\item The Riemann Hypothesis holds.
\item $Q(\Phi) \geq 0$ for every $\Phi \in \mathcal{W}$.
\end{enumerate}
\end{theorem}

\begin{proof}[Proof sketch]
The equivalence follows from the Guinand--Weil explicit formula, which connects the zeros of $\zeta(s)$ to the functional $Q$. If RH holds, then all non-trivial zeros contribute positively to a certain representation of $Q$, yielding (ii). Conversely, if (ii) holds and there were a zero off the critical line, one could construct a test function making $Q$ negative, contradicting (ii). The full proof appears in Weil~\cite{Weil1952,Weil1972}; see also~\cite[Ch.~5]{Edwards1974} and~\cite[Ch.~5]{IwaniecKowalski2004}.
\end{proof}

\begin{remark}[Explicit formula tradition]\label{rem:explicit}
Weil's criterion builds on the explicit formula tradition: Riemann's original formula~\cite{Riemann1859} connecting primes to zeros, Guinand's summation formula~\cite{Guinand1948}, and Weil's reformulation as a positivity statement~\cite{Weil1952}. The criterion reduces the infinitely many constraints (one for each zero) to a single positivity condition.
\end{remark}

%=============================================================================
\section{Proof of the Main Theorem}
\label{sec:RH}
%=============================================================================

We now complete the proof of the Riemann Hypothesis.

\begin{theorem}[Riemann Hypothesis]\label{thm:RH}
All non-trivial zeros of the Riemann zeta function $\zeta(s)$ lie on the critical line $\Re(s) = 1/2$.
\end{theorem}

\begin{proof}
By Theorem~\ref{thm:main-pos} (Main positivity), we have $Q(\Phi) \geq 0$ for all $\Phi$ in the Weil cone $\mathcal{W}$. By Theorem~\ref{thm:weil-criterion} (Weil's criterion), this implies the Riemann Hypothesis.
\end{proof}

\begin{remark}[Scope of the proof]\label{rem:scope}
The proof establishes RH under the analytic module stack (T0) + (A1') + (A2) + (A3) + (RKHS). Each module is proved by explicit analytic estimates in Parts~I, II, and this paper. No numerical computations, computer algebra, or automated theorem provers are used in the main argument.
\end{remark}

\begin{remark}[The module architecture]\label{rem:architecture}
The modular structure allows each component to be verified independently:
\begin{itemize}
\item (T0) is a straightforward coordinate change.
\item (A1') is classical approximation theory.
\item (A2) is Lipschitz analysis with explicit constants.
\item (A3) combines Szeg\H{o}--B\"ottcher theory with symbol analysis.
\item (RKHS) exploits the Gram matrix structure of the heat kernel.
\end{itemize}
The synthesis in Theorem~\ref{thm:main-pos} is a short argument combining these inputs.
\end{remark}

%=============================================================================
\section{Discussion}
\label{sec:discussion}
%=============================================================================

\subsection{Summary of the trilogy}

This trilogy develops an operator-theoretic proof of the Riemann Hypothesis:

\begin{itemize}
\item \textbf{Part~I}~\cite{PartI} establishes the foundational properties (A1') and (A2): density of the Fej\'er--heat cone and Lipschitz continuity of the Weil functional.

\item \textbf{Part~II}~\cite{PartII} establishes the spectral analysis (A3) and (RKHS): the Toeplitz barrier via Szeg\H{o}--B\"ottcher asymptotics and the prime cap via RKHS contraction.

\item \textbf{Part~III} (this paper) synthesizes the modules to prove $Q \geq 0$ on the Weil cone, and applies Weil's criterion to establish RH.
\end{itemize}

\subsection{Key constants}

The proof involves the following explicit constants:

\begin{center}
\begin{tabular}{@{}lll@{}}
\toprule
Constant & Value & Role \\
\midrule
$c_*$ & $11/10$ & Archimedean floor \\
$C_{\mathrm{SB}}$ & $4$ & Szeg\H{o}--B\"ottcher constant \\
$t_{\mathrm{sym}}$ & $3/50$ & Symbol heat parameter \\
$t_{\mathrm{rkhs}}$ & $\geq 1$ & RKHS threshold \\
$B_{\min}$ & $3$ & Minimum bandwidth \\
$w_{\max}$ & $\leq 2/e$ & Maximum prime weight \\
$c_*/4$ & $11/40$ & Spectral gap \\
\bottomrule
\end{tabular}
\end{center}

\subsection{Verification philosophy}

The proof is designed for verifiability:
\begin{enumerate}
\item All bounds are explicit analytic estimates.
\item No $K$-dependent schedules or optimization is required.
\item Each module can be checked independently.
\item The synthesis (Theorem~\ref{thm:main-pos}) is a short, transparent argument.
\end{enumerate}

\subsection{Open questions}

Several natural extensions arise:

\begin{enumerate}
\item \textbf{Generalized Riemann Hypothesis:} Can the module chain be extended to Dirichlet $L$-functions?

\item \textbf{Explicit zero-free regions:} Can the spectral gap $c_*/4$ be converted to an explicit zero-free region?

\item \textbf{Simplification:} Can the module chain be shortened or the constants improved?

\item \textbf{Alternative architectures:} Are there other module decompositions that yield RH?
\end{enumerate}

%=============================================================================
% Bibliography
%=============================================================================

\begin{thebibliography}{99}

\bibitem{PartI}
Authors,
``Fej\'er--heat generators and Lipschitz control for the Weil quadratic functional,''
Part~I of this series.

\bibitem{PartII}
Authors,
``Toeplitz barrier and RKHS prime contraction for the Weil positivity criterion,''
Part~II of this series.

\bibitem{Edwards1974}
H.~M.~Edwards,
\emph{Riemann's Zeta Function},
Academic Press, New York, 1974.

\bibitem{GriffinOnoRolenZagier2019}
M.~Griffin, K.~Ono, L.~Rolen, and D.~Zagier,
``Jensen polynomials for the Riemann zeta function and other sequences,''
\emph{Proc.\ Natl.\ Acad.\ Sci.\ USA}\ \textbf{116} (2019), 11103--11110.

\bibitem{Guinand1948}
A.~P.~Guinand,
``A summation formula in the theory of prime numbers,''
\emph{Proc.\ London Math.\ Soc.}\ \textbf{50} (1948), 107--119.

\bibitem{IwaniecKowalski2004}
H.~Iwaniec and E.~Kowalski,
\emph{Analytic Number Theory},
AMS Colloquium Publications, Vol.~53, 2004.

\bibitem{Li1997}
X.-J.~Li,
``The positivity of a sequence of numbers and the Riemann hypothesis,''
\emph{J.\ Number Theory}\ \textbf{65} (1997), 325--333.

\bibitem{Riemann1859}
B.~Riemann,
``\"Uber die Anzahl der Primzahlen unter einer gegebenen Gr\"osse,''
\emph{Monatsberichte der Berliner Akademie} (1859), 671--680.

\bibitem{Titchmarsh1986}
E.~C.~Titchmarsh,
\emph{The Theory of the Riemann Zeta-Function},
2nd ed.\ (revised by D.~R.~Heath-Brown), Oxford University Press, Oxford, 1986.

\bibitem{Weil1952}
A.~Weil,
``Sur les 'formules explicites' de la th\'eorie des nombres premiers,''
\emph{Comm.\ S\'em.\ Math.\ Univ.\ Lund, Tome Suppl\'ementaire} (1952), 252--265.

\bibitem{Weil1972}
A.~Weil,
``Sur les formules explicites de la th\'eorie des nombres,''
\emph{Izv.\ Akad.\ Nauk SSSR Ser.\ Mat.}\ \textbf{36} (1972), 3--18.

\end{thebibliography}

\end{document}
